\documentclass[a4paper]{article}
\usepackage{MMVII}

\newcommand{\TheCmd}{OriPoseEstimRel2Im}

\begin{document}

\CmdTitle
%\maketitle

\tableofcontents


%=================================================================================

\SecDescr

This command  compute the relative orientation of a pair of images
using tie points and a known calibration for each images.
The main purpose of this functionnality it to be used
when computing  global orientation   by  other commands :

\begin{itemize}
    \item \FileSty{OriPoseEstimRelPairsOf1Im} 
    \item \FileSty{OriPoseEstimRelAllPairs}
\end{itemize}

Having access to thi elementary command is (almost) never done when
computing a real 3D modelization.
However, it  may be interesting to run on a specific pair for :


\begin{itemize}
   \item didactic purpose when teaching the methodology for computing global
         relative orientation;

   \item tuning and testing when developping the complete pipeline.

\end{itemize}


Several parameters used for  this command will be used similarly for  other
command on pair and triplets (see~\ref{Ref:RelCmd}).

The $4$ mandatory parameters are:

\begin{itemize}
  \item the $2$ images we want to compute relative orientation ( pattern, or triplet for other command);

  \item the input folder for localization of internal calibration of the $2$ images (similar in command of~\ref{Ref:RelCmd});

  \item an output folder for saving the $2$ relative poses as a standard orientation
       (because of the "didactic" , for other command, it will be saved in folder specific 
        aspect).
\end{itemize}


%=================================================================================

\section{"Mandatory" Optionnal parameters for input}

All these command accept $3$ different way to specify the tie points used to compute orientation :

\begin{itemize}
   \item  tie point computed from pair of images with parameter \FileSty{InTieP};
	  this is the \emph{standar} method;

   \item  multiple tie point issued from merged tie points by pair , with
          parameter \FileSty{InMulTieP};
	  this can be  used when point by pair have been merged, and initial 
	  by pair have been "lost";

  \item  points computed from coded target with parameter \FileSty{InObjMesInstr},
	 this can be used when computing  relative orientation with code-target
         that have no $3d$ measure associated.
\end{itemize}

%=================================================================================

\section{Methodology used}

    %  - - - - - - - - - - - - - - - - - - - - - - - - - - - - - - - - - - - - - - - - - - - - - - - - - - - -

\subsection{Introduction}

We give a very brief description of the method used :

\begin{itemize}
    \item we test different algorithm for initial estimation;
    \item for each solution we make refinement by bundle adjustment
    \item we select the best refined solution  using a criteria on
	    reprojection that try to be robust to outliers.
\end{itemize}

    %  - - - - - - - - - - - - - - - - - - - - - - - - - - - - - - - - - - - - - - - - - - - - - - - - - - - -

\subsection{Initialization}

We test $2$ method :

\begin{itemize}
    \item  the essential matrix that works in all case, except when the scene is planar
    \item  the homographic decomposition that works only when the scene is planar;
	   is the scene is planar, the method gives $2$ solution that cannot be distinguished;
\end{itemize}

For each method we make ransac estimation.
So we obtain $3$ possible initial solutions that are refined using a "specialized bundled adjustment".

    %  - - - - - - - - - - - - - - - - - - - - - - - - - - - - - - - - - - - - - - - - - - - - - - - - - - - -

\subsection{Selection}

Let $r_0,r_1, \dots r_{N-1}$ be the $N$ residual of the reprojection of tie points,
\emph{sorted} in growing order,
the formula used for comparing the different solution is : 

\begin{equation}
	\sum_{k=0}^N  r_k (N-k)
\end{equation}



%=================================================================================

\section{Parameters for testing and tuning}

\section{Ground truth}

 %=================================================================================
\SecRelatedCmd

\begin{itemize}
   \item  \FileSty{OriPoseEstimRelPairsOf1Im}
   \item  \FileSty{OriPoseEstimRelAllPairs}
   \item  \FileSty{OriPoseEstimRel3Im}
   \item  \FileSty{OriPoseEstimRelTripletssOf1Im}
   \item  \FileSty{OriPoseEstimRelAllTriplets}
   \item  \FileSty{OriPoseEstimRelGlob}
\end{itemize}

And 


\end{document}
